% ============================================================
%  SECCIÓN — PROCESO CONSTRUCTIVO (MURO)
% ============================================================
\section{Proceso constructivo}
\label{sec:muro_proceso}

% -----------------------------------------------------------
\subsection{Descripción general}
% -----------------------------------------------------------

El muro de contención en ménsula de $L = 150$~ml y desnivel $H = 7{,}30$~m se
ejecuta mediante un procedimiento de \textbf{obra nueva en desmonte abierto},
trabajando desde el lado libre (intradós) sin necesidad de contención provisional
del terreno. La cimentación es mixta: zapata superficial de HA-30 de
$B = 5{,}60$~m reforzada con pilotes CPI ($D = 0{,}45$~m) debido al asiento
excesivo estimado ($s = 7{,}51$~cm $>$ 5~cm admisible).

% -----------------------------------------------------------
\subsection{Descripción paso a paso}
\label{subsec:proceso_pasos}
% -----------------------------------------------------------

\subsubsection*{1. Replanteo y trabajos previos}

Se realiza el replanteo topográfico del eje del muro con estación total, marcando
los límites de la zapata con estacas y cordel de referencia. Se señaliza el
perímetro de obra y se garantiza el acceso de maquinaria por el lado libre
(intradós).

\subsubsection*{2. Excavación de la zanja de cimentación}

Se ejecuta una \textbf{excavación a cielo abierto} para la implantación de la
zapata desde el intradós. La plataforma de trabajo tiene un ancho mínimo:
\[
  B_{\mathrm{exc}} = B + 2 \times 0{,}60 = 5{,}60 + 1{,}20
  = \mathbf{6{,}80~\text{m}}
\]
para permitir la maniobra de encofrado. La profundidad de excavación desde la
rasante del intradós es $D_f = 1{,}00$~m (cota de base de zapata).
El material excavado se transporta a vertedero autorizado.

\begin{notabox}[Condicionantes medioambientales]
  Dado el carácter medioambientalmente protegido de la zona y la presencia de
  sulfatos ($3\,500$~mg~SO${}_4^{2-}$/kg), la tierra excavada se gestiona como
  residuo de tierras según la normativa vigente. Se evitará el contacto de
  lodos o aguas con sulfatos con el hormigón fresco.
\end{notabox}

\subsubsection*{3. Capa de hormigón de limpieza (HM-15)}

Sobre el fondo de la zanja se extiende y nivela una capa de
\textbf{hormigón de limpieza HM-15} de $e = 10$~cm de espesor. Proporciona
una superficie limpia y uniforme para la ferralla y evita la contaminación del
acero por el terreno.
\[
  V_{\mathrm{HM-15}} = 5{,}60 \times 0{,}10 \times 150 = \mathbf{84~\text{m}^3}
\]

\subsubsection*{4. Pilotaje CPI ($D = 0{,}45$~m)}

Dado que el asiento estimado de la zapata superficial es de $s = 7{,}51$~cm
$> 5$~cm admisible, se ejecuta \textbf{cimentación profunda mediante pilotes
CPI} (barrena helicoidal continua) de $D = 0{,}45$~m, a razón de un pilote
por metro lineal de muro (150 pilotes en total).

El proceso de ejecución de cada pilote es:
\begin{enumerate}
  \item Posicionamiento de la perforadora sobre el eje del pilote según replanteo.
  \item Perforación con barrena helicoidal continua hasta empotrar $0{,}5$~m
        en el estrato de arcillas margosas (cota de base $\approx -8{,}6$~m desde
        trasdós), que es el estrato resistente de cimentación.
  \item Inyección de hormigón HA-30 autocompactante por el núcleo hueco de
        la barrena \emph{durante la extracción}, evitando el colapso de paredes.
  \item Colocación de la armadura de espera por vibrado antes del inicio del
        fraguado.
\end{enumerate}

\subsubsection*{5. Ferrallado de la zapata}

Una vez endurecido el hormigón de limpieza, se coloca la
\textbf{armadura de la zapata} en el siguiente orden:
\begin{enumerate}
  \item Distanciadores para garantizar el recubrimiento nominal
        $r_{\mathrm{nom}} = 50$~mm (clase de exposición IIb+Qc).
  \item Armadura superior del talón (cara superior): $6\,\phi 16$~c/16{,}7~cm.
  \item Armadura inferior de la puntera (cara inferior): $6\,\phi 16$~c/16{,}7~cm.
  \item Arranques del alzado: barras verticales $\phi 16$ con solape mínimo
        según el Código Estructural (CE~2021).
\end{enumerate}

\subsubsection*{6. Encofrado y hormigonado de la zapata}

El encofrado de la zapata se realiza con \textbf{encofrado metálico} a ambos
lados, con testeros en las juntas de dilatación cada 15~m.
Se hormigona con HA-30/B/20/IIb+Qc (cemento SR) mediante bomba estacionaria,
vibrado con aguja eléctrica.
\[
  V_{\mathrm{zapata}} = B \times h_z \times L
  = 5{,}60 \times 0{,}80 \times 150 = \mathbf{672~\text{m}^3}
\]

\subsubsection*{7. Curado de la zapata}

El hormigón de la zapata se cura durante un \textbf{mínimo de 7 días}
manteniéndolo húmedo o con producto de curado en las superficies expuestas.
El desencofrado no se realizará antes de que el hormigón alcance el 70~\% de
su resistencia característica.

\subsubsection*{8. Ferrallado del alzado}

Con la zapata curada se procede al ferrallado del \textbf{alzado}:
\begin{enumerate}
  \item Empalme por solape de los arranques con la armadura principal del
        trasdós: $6\,\phi 16$~c/16{,}7~cm (armadura principal, cara trasdós).
  \item Estribos de cortante: 3~ramas $\phi 10$~c/19{,}5~cm.
  \item Separadores y galgas de encofrado.
\end{enumerate}

\subsubsection*{9. Encofrado del alzado}

Se adopta \textbf{encofrado metálico a dos caras}, con altura máxima de panel
de $\approx 3{,}25$~m, ejecutado en dos tongadas de altura. El alzado tiene
sección variable ($e_1 = 0{,}60$~m en coronación, $e_2 = 0{,}75$~m en la
base), lo que requiere ajustar los separadores en cada tongada.
\[
  S_{\mathrm{encofrado}} = 2 \times H \times L
  = 2 \times 7{,}30 \times 150 = \mathbf{2\,190~\text{m}^2}
\]

\subsubsection*{10. Hormigonado del alzado}

El hormigonado del alzado se realiza en \textbf{tongadas de altura controlada}
para limitar la presión sobre el encofrado. Se vierte con bomba estacionaria
y se vibra cada tongada antes de continuar con la siguiente.
\[
  V_{\mathrm{fuste}} = \frac{e_1+e_2}{2} \times H \times L
  = \frac{0{,}60+0{,}75}{2} \times 7{,}30 \times 150
  = \mathbf{740~\text{m}^3}
\]

\subsubsection*{11. Desencofrado y curado del alzado}

Desencofrado no antes de las \textbf{72~h} desde el hormigonado de cada
tongada. Curado mínimo de 7~días. Sellado de marcas de tirantes con mortero
de reparación de igual clase de exposición.

\subsubsection*{12. Impermeabilización del trasdós}

Para proteger el muro del ataque de sulfatos del suelo y garantizar la
estanqueidad, se aplica un sistema de impermeabilización continuo sobre el
trasdós del alzado:
\begin{enumerate}
  \item \textbf{Imprimación bituminosa} sobre el hormigón limpio y seco.
  \item \textbf{Lámina impermeabilizante de polietileno de alta densidad (HDPE)}
        de $e \geq 2$~mm, fijada mecánicamente en la coronación y sellada en
        los solapes.
  \item \textbf{Lámina drenante nodular} para conducir el agua infiltrada
        hacia el tubo de drenaje de la base.
  \item \textbf{Geotextil de protección} sobre la lámina drenante para evitar
        la colmatación.
\end{enumerate}

\subsubsection*{13. Instalación del sistema de drenaje}

En la base del trasdós, sobre el canto superior de la zapata y antes del
relleno, se coloca un \textbf{tubo de drenaje perforado de $\varnothing 100$~mm}
rodeado de grava lavada filtrada con geotextil. El tubo desagua cada 15~m a
arquetas de recogida conectadas a la red de pluviales.

\subsubsection*{14. Relleno del trasdós por tongadas}

El relleno del trasdós se ejecuta por \textbf{tongadas horizontales de 30~cm
de espesor}, compactado con rodillo ligero o placa vibratoria hasta alcanzar
el \textbf{95~\% del Proctor Modificado}. No se utilizarán compactadores
pesados junto al trasdós para no generar sobrepresiones sobre el muro.

\begin{notabox}[Material de relleno]
  Se utilizará zahorra natural o material granular de baja cohesión y buena
  permeabilidad, compatible con el sistema de drenaje instalado.
\end{notabox}

\subsubsection*{15. Juntas de dilatación y acabados}

Se ejecutan \textbf{juntas de dilatación cada 15--20~m} mediante junta de
poliestireno expandido de 20~mm sellada con masilla de poliuretano.
Remate de coronación del muro y reposición del firme del camino en coronación
(zahorra compactada o mezcla bituminosa según el acceso existente).

% -----------------------------------------------------------
\subsection{Medios humanos}
\label{subsec:medios_humanos}
% -----------------------------------------------------------

\begin{table}[H]
  \centering
  \caption{Cuadrillas y equipos previstos para las unidades principales}
  \label{tab:cuadrillas}
  \resizebox{\textwidth}{!}{%
  \begin{tabular}{@{}llcc@{}}
    \toprule
    \textbf{Unidad de obra} & \textbf{Cuadrilla / Equipo}
    & \textbf{Rendimiento} & \textbf{Duración aprox.} \\
    \midrule
    Excavación a máquina  & 2 peones + retroexcavadora & 300~m³/día  & 7~días \\
    HM-15 de limpieza     & 1 oficial + 1 peón + bomba & 30~m³/día   & 3~días \\
    Pilotaje CPI          & Equipo especializado       & 10~pilotes/día & 15~días \\
    Ferrallado            & 4 oficiales ferralistas    & 1\,000~kg/día & 45~días \\
    Encofrado metálico    & 2 oficiales + 2 peones     & 40~m²/día   & 55~días \\
    Hormigonado (bomba)   & 2 oficiales + 2 peones     & 30~m³/h     & 45~h \\
    Curado                & 1 peón                     & —           & 7~días/fase \\
    Impermeabilización    & Cuadrilla especializada    & 200~m²/día  & 11~días \\
    Relleno y compactación & 2 peones + rodillo         & 150~m³/día  & 22~días \\
    \bottomrule
  \end{tabular}}
\end{table}

% -----------------------------------------------------------
\subsection{Maquinaria y medios auxiliares}
\label{subsec:maquinaria}
% -----------------------------------------------------------

Los \textbf{medios auxiliares} son equipos o instalaciones de carácter
provisional que facilitan la ejecución de la obra pero no forman parte de
la estructura definitiva. En el presente proyecto los más relevantes son los
descritos a continuación.

\paragraph{Grúa automontante (tipo torre).}
Máquina de elevación fija formada por mástil vertical, pluma horizontal y
carro desplazador de la carga. Se emplea para el acopio y montaje de
ferralla y paneles de encofrado a lo largo de los $150$~ml de muro. La
selección se realiza mediante el \emph{diagrama de cargas} del fabricante,
que relaciona la carga máxima levantable con la distancia al eje de la grúa:
se adoptan unidades con pluma de $40$~m y carga útil $\geq 1\,500$~kg a
$30$~m de radio, cubiertas por \textbf{plan de maniobras} obligatorio
(cargas a elevar, ubicación, coeficientes de seguridad, elementos de izado).

\paragraph{Entibación de la zanja de cimentación.}
La zanja tiene una profundidad de $D_f = 1{,}00$~m, con el nivel freático
a $-2{,}50$~m, por lo que el riesgo de colapso de paredes es moderado. Se
prevé \textbf{entibación ligera} (módulos de aluminio, fáciles de montar)
de tipo \emph{semicuajada}, que cubre aproximadamente el 50~\% de la
superficie de cada pared lateral. Alternativamente puede evitarse
entibando si se talúa el terreno o se ejecuta la prezanja con mayor anchura
y menor profundidad, tal como contempla el proceso descrito en el paso~2.

\paragraph{Elevadores hidráulicos.}
Para los trabajos en altura sobre el alzado ($H = 7{,}30$~m) —encofrado,
cosido de estribos, impermeabilización del trasdós— se utilizan
\textbf{elevadores telescópicos/articulados} con brazo articulado de
alcance horizontal hasta $38$~m, capaces de salvar obstáculos como las
armaduras ya colocadas. Este tipo de plataforma es preferible a los
elevadores de tijera (limitados a $8$–$20$~m en vertical estricto) por la
necesidad de alcance lateral sobre el trasdós durante la aplicación de
la lámina HDPE.

\paragraph{Bomba sumergible.}
Con el nivel freático a $-2{,}50$~m y la base de la zanja a $-1{,}00$~m
la excavación puede permanecer seca; no obstante, lluvias intensas o
variaciones estacionales del acuífero obligan a disponer una
\textbf{bomba sumergible} (motor + bomba + tubería de impulsión + filtro
\emph{alcachofa}) para el achique rápido de la zanja, el drenaje de
posibles bolsas de agua durante el pilotaje y el vaciado de la plataforma
de trabajo.

\paragraph{Grupo electrógeno.}
Dado el carácter lineal de la obra ($150$~ml en zona periurbana), no está
garantizado el suministro eléctrico permanente en todos los frentes de
trabajo. Se dispone de un \textbf{grupo electrógeno} portátil para
alimentar la aguja de vibrado, la bomba sumergible y el alumbrado
provisional de obra.

\paragraph{Compresor.}
Genera aire comprimido para los \textbf{martillos neumáticos} y equipos de
perforación empleados en el ajuste del encofrado y en las operaciones de
limpieza superficial del hormigón.

\begin{table}[H]
  \centering
  \caption{Maquinaria principal prevista}
  \label{tab:maquinaria}
  \resizebox{\textwidth}{!}{%
  \begin{tabular}{@{}lll@{}}
    \toprule
    \textbf{Máquina} & \textbf{Características} & \textbf{Uso} \\
    \midrule
    Retroexcavadora mixta
      & Capacidad $\approx 0{,}80$~m³ — 2~unidades
      & Excavación y carga \\
    Perforadora CPI (barrena helicoidal)
      & $D_{\max} \geq 0{,}60$~m
      & Pilotaje $D = 0{,}45$~m \\
    Camión dumper
      & 20~Tn — 3~unidades
      & Transporte a vertedero \\
    Camión hormigonera
      & 8~m³ — suministro de HA-30
      & Hormigonado \\
    Bomba estacionaria
      & Alcance $\geq 40$~m
      & Zapata y alzado \\
    Grúa automontante
      & Alcance 40~m, carga útil $\geq 1\,500$~kg a 30~m
      & Ferralla y encofrado \\
    Rodillo vibrante ligero
      & $\leq 1{,}0$~Tn
      & Compactación trasdós \\
    Placa vibratoria
      & $\approx 500$~kg
      & Trasdós junto al muro \\
    \bottomrule
  \end{tabular}}
\end{table}

\paragraph{Retroexcavadora mixta.}
Se utilizan \textbf{2~unidades} para abordar los $1\,938$~m³ de excavación en
terreno cohesivo (arcillas limosas) con continuidad a lo largo de los $150$~ml.
La configuración mixta —brazo retroexcavador trasero más pala frontal— permite
excavar la zanja de $6{,}80 \times 1{,}90$~m y cargar los camiones dumper con una
sola máquina; las orugas garantizan tracción y estabilidad en el fondo de
excavación sin necesidad de banqueta.

\paragraph{Perforadora CPI con barrena helicoidal continua.}
La necesidad de ejecutar $150$~pilotes de $D = 0{,}45$~m hasta $-8{,}00$~m en arcillas
con nivel freático a $-2{,}50$~m hace que la técnica CPI con barrena de eje hueco sea
la más apropiada: el hormigón se bombea por el interior de la barrena durante su
extracción, de modo que la perforación queda rellena antes de retirarla del
terreno, eliminando el riesgo de sifonamiento y la necesidad de entubación o
lodos bentoníticos (incompatibles con la zona protegida).

\paragraph{Camión dumper.}
Se prevén \textbf{3~unidades} de $20$~Tn para gestionar el flujo de material
excavado sin crear cuellos de botella frente a las dos retroexcavadoras. Su caja
basculante permite la descarga directa en vertedero autorizado y su tracción
todo-terreno cubre la rampa provisional de acceso a la zanja.

\paragraph{Camión hormigonera.}
El volumen total de HA-30/B/20/IIb+Qc asciende a $1\,331$~m³, lo que exige
suministro continuo desde planta externa. La cuba de $8$~m³ por camión,
coordinada con la bomba estacionaria, garantiza que no se produzcan juntas frías
durante las coladas de zapata ($672$~m³) y alzado ($659$~m³).

\paragraph{Bomba estacionaria de hormigón.}
El alzado del muro alcanza $6{,}50$~m de altura y la zapata tiene $5{,}60$~m de
anchura, por lo que el camión hormigonera no puede verter directamente en el
encofrado. La bomba, con alcance $\geq 40$~m, se posiciona en el coronamiento de
la excavación y llega a cualquier punto de la sección con un solo tendido de
tubería, limitando los reposicionamientos a uno cada $\approx 40$~ml.

\paragraph{Grúa automontante.}
La variante \textbf{automontante} se erige por sus propios medios sin necesitar
grúa auxiliar, lo que resulta idóneo en una obra lineal de $150$~ml con acceso
limitado. Con un alcance de $40$~m y carga útil $\geq 1\,500$~kg a $30$~m de radio,
cubre el izado de las jaulas de armadura del alzado, los paneles metálicos de
encofrado y la descarga de la perforadora CPI a lo largo de toda la alineación.
La selección del modelo se verifica mediante su diagrama de cargas en la posición
más desfavorable (mayor radio y carga máxima).

\paragraph{Rodillo vibrante ligero.}
El trasdós del muro presenta un espacio estrecho entre la cara del fuste y el
talud de la excavación, suficiente para un rodillo de lanza ($\leq 1{,}0$~Tn) pero
no para maquinaria pesada. Este equipo compacta la zahorra en capas de
$20{-}30$~cm con mayor energía que una placa vibratoria, reduciendo el número de
pasadas necesarias para los $3\,949$~m³ de relleno previstos.

\paragraph{Placa vibratoria.}
Se emplea exclusivamente en las zonas donde el rodillo de lanza no puede maniobrar:
junto al intradós del muro y en las esquinas de la zapata. Trabaja en capas finas
($\leq 15$~cm) y es el único equipo capaz de compactar esos rincones sin transmitir
empujes laterales que puedan dañar el encofrado o el hormigón fresco del alzado.
